% ============================================================
% BME VIK TDK abstract — VARIANT 3: novelty-as-residual (F4)
%
% Third of the framings for the supervisors to compare. Variant 1
% (abstract-v1-problem-method.tex) frames the work as reconstructing
% an abstract from the citation graph (capability claim). Variant 2
% (abstract-v2-explainability.tex) frames the same system as a
% diagnostic testbed for graph-conditioned LLMs. This variant keeps
% the same machinery and inverts the objective: the part of a paper
% its bibliography FAILS to predict is the measurement, not the error.
%
% Why this framing: reconstruction from a bibliography has a hard
% information ceiling, because a paper's contribution is by definition
% its departure from the work it cites. V1 bets against that ceiling.
% This variant measures it.
%
% Structure follows the official one-page sample:
% https://portal.vik.bme.hu/files/00000047.pdf
%   - max 1 page, min ~20 lines of body text
%   - NO header, footer, or page numbers
%   - figures/equations only if strictly necessary
%   - own contribution must be clearly separated from cited work
%   - body must cover: methodology, key results, conclusions
%
% NUMBERS: the body is written so that measured figures slot into the
% validation paragraph (rank correlation against the disruption index,
% landmark-vs-control separation, calibration accuracy) without
% restructuring. Do not submit without them — the official sample
% requires key results in the body.
%
% NOTE ON LANGUAGE: the official TDK conference language is Hungarian
% (English talks require prior agreement with the section chair, per
% the official sample). This English draft is for internal review with
% Bálint/Dani; translate to Hungarian before submission unless an
% English exception is confirmed.
% ============================================================

\documentclass[12pt]{article}

\usepackage[utf8]{inputenc}
\usepackage[T1]{fontenc}
\usepackage[english]{babel}
\usepackage[a4paper, margin=2.5cm]{geometry}
\usepackage{mathptmx}      % Times-like font
\usepackage{hyperref}
\usepackage{enumitem}

\pagestyle{empty}          % <- enforces "no header/footer/page numbers"
\setlength{\parindent}{0pt}
\setlength{\parskip}{0.6em}

\begin{document}

% ---------- Title ----------
\begin{center}
    {\Large\bfseries What a Bibliography Cannot Predict: \\
    Scientific Novelty as a Generative Residual}
\end{center}

\vspace{0.8em}

% ---------- Authors ----------
\begin{center}
    Gema Csanád IV. Inf, gemacsanad@gmail.com \\

\end{center}

\vspace{0.4em}

% ---------- Consultant ----------
\begin{center}
    Consultant: Unyi Dániel
\end{center}

\vspace{1.2em}

% ==========================================================
% BODY
% ==========================================================

A citation graph represents scientific literature as a network: each
node is a publication, each edge a citation from one paper to
another, and every node carries that paper's abstract as its text
content. A reference list is not an arbitrary collection — it records
the problem a paper attacks, the methods it borrows and the community
it belongs to. This work asks how much of a paper that record
determines, and treats what it does not determine as a quantity that
can be measured.

A paper's contribution is by definition its departure from the work
it cites: anything already contained in its references would not be a
contribution. We turn this into a measurement. A model is trained to
predict a paper's abstract from its citation neighbourhood alone,
with the paper's own text withheld entirely, and the discrepancy
between the predicted and the actual abstract is taken as an estimate
of the paper's novelty. What makes reconstruction from a bibliography
impossible in principle is exactly what this measure captures.

Established measures approach novelty from two directions, neither
conditioned on an individual paper's own antecedents. Reference-based
measures score how atypical the combination of cited journals is,
discarding content altogether. Language-model measures score the
perplexity of a paper's text under a general-purpose model, asking
whether the text is unusual in general rather than unusual given what
this paper was built on. A third line has a language model read a
paper and judge its novelty, which inherits the judge's priors rather
than measuring anything. Our measure is conditional and generative:
the expectation a paper is scored against is built from its own
reference list.

Methodologically, a graph neural network encodes the citation
neighbourhood — the cited papers, and their references in turn,
beyond what could fit in a prompt — into a continuous representation,
which a trained adapter projects into soft prefix tokens injected
into the attention layers of a language model. Generation is
conditioned on this representation together with the text of the
directly cited papers. The target's own features are masked, and
edges are oriented from cited to citing paper, so nothing published
after the target can reach it. The residual is scored separately for
research domain, problem, task and methodology, because a
bibliography determines these to very different degrees.

Such a residual confounds genuine novelty with model error and with
incomplete bibliographic coverage. We address this in two ways: the
measure is relative rather than absolute, read against the residual
distribution of matched papers of the same field, year and observed
reference count; and the model is first calibrated on papers whose
content the bibliography does determine, since a residual is evidence
of novelty only insofar as the model is accurate where prediction is
possible.

Our own contribution is the formulation of novelty as a conditional
generative residual, the graph-conditioned model producing it, and
its validation — the substantive difficulty, since novelty has no
ground truth. We validate without one: against axiomatic properties
any novelty measure should satisfy, by separation of recognised
landmark papers from matched controls, and by correlation with the
disruption index and later citation impact computed on the same
graph. We report where the measure agrees with existing reference-
and perplexity-based scores and, more informatively, where it
disagrees.

\end{document}
